\color{unidarkgreen}{\subsection{Токовая направленная защита нулевой последовательности (ТНЗНП)}}
\color{black}

\begin{enumerate}[label=\arabic{section}.\arabic{subsection}.\arabic*,  align=left, labelsep=4pt, leftmargin=0pt, itemindent=57pt]

\item
Токовая направленная защита нулевой последовательности (ТНЗНП) предназначена для селективного
отключения поврежденного участка электрической сети при коротких замыканиях на землю, реагируя на
повышение тока нулевой последовательности.

\item
В составе логики токовой направленной защиты нулевой последовательности предусмотрено четыре
ступени от замыканий на землю.

\vspace{2mm}
\begin{figure}[H]
\centering
\includegraphics[width=1\textwidth,height=1\textheight,keepaspectratio]{\fbpath/661.0121 ТНЗНП/img/TNZNP.pdf}
\captionof{figure}{Функциональная схема ТНЗНП-1(2…4}\label{pict_tnznp:TNZNP}
\end{figure}

\item
Пороги тока срабатывания ступеней ТНЗНП задаются уставками \textbf{«3I0ср ТНЗНП-1(2…4)»}. Действие
ступеней осуществляется с регулируемыми независимыми выдержками времени \textbf{«Тср ТНЗНП-1(2…4)»}. Ввод
в работу ступеней ТНЗНП осуществляется с помощью программных ключей \textbf{«ТНЗНП-1(2…4)»}. 

\item
Направленность ТНЗНП обеспечивается органами направления мощности нулевой последовательности
(ОНМНП) и обратной последовательности (ОНМОП).

\item
ОНМНП реагирует на величины векторов тока и напряжения нулевой последовательности от ТН 1СШ, а
также угол сдвига между ними.

\item
Так как мощность нулевой последовательности всегда направлена от места несимметрии во внешнюю
сеть, то ОНМНП в 1СШ срабатывает при направлении мощности нулевой последовательности от 1СШ к ШСВ
(при КЗ впереди), а  ОНМНП в сеть – при обратном направлении мощности. 

\item
За положительное направление отсчета векторов токов и напряжений принято направление против
часовой стрелки. Углы максимальной чувствительности откладываются от вектора напряжения 3U0 (вектор
поляризации) и задаются уставками \textbf{«Фмч ОНМНП в 1СШ»} и \textbf{«Фмч ОНМНП в 2СШ»}. Уставки по углам
максимальной чувствительности рекомендуется задавать для ОНМНП в 1СШ – плюс 110º, а для ОНМНП в
2СШ – плюс 290º. Ширина зоны срабатывания ± 85º от угла максимальной чувствительности.

\item
Пороги срабатывания ОНМНП в 1СШ и ОНМНП в 2СШ по току и напряжению регулируются уставками 
\textbf{«3I0ср ОНМНП в 1СШ»}, \textbf{«3U0ср ОНМНП а 1СШ»}, \textbf{«3I0ср ОНМНП в 2СШ»} и \textbf{«3U0ср ОНМНП в 2СШ»}. Выбор
напряжения нулевой последовательности, по которому работает ОНМНП, осуществляется с помощью
программного переключателя \textbf{«Выбор 3U0»}. 

\item
Орган направления мощности обратной последовательности ОНМОП в 1СШ срабатывает при
направлении мощности обратной последовательности от 1СШ к ШСВ, а ОНМОП в 2СШ -- при обратном
направлении мощности. 

\item
Углы максимальной чувствительности задаются уставками \textbf{«Фмч ОНМОП в 1СШ»} и \textbf{«Фмч ОНМОП в
2СШ»}. Уставки по углам максимальной чувствительности рекомендуется задавать для ОНМОП в 1СШ -- плюс
$110^\circ$, а для ОНМОП в 2СШ -- плюс $290^\circ$. Ширина зоны срабатывания $\pm 85^\circ$ от угла максимальной
чувствительности.

\item
Пороги срабатывания ОНМОП в 1СШ и ОНМОП в 2СШ  по току и напряжению регулируются уставками
\textbf{«3I0ср ОНМОП в 1СШ»}, \textbf{«3U0ср ОНМОП в 1СШ»}, \textbf{«3I0ср ОНМОП в 2СШ»}, и \textbf{«3U0ср ОНМОП в 2СШ»}.

\small
\setlength{\extrarowheight}{0.05cm} % высота ячеек в таблице
\begin{longtable}
{|>{\centering\arraybackslash}m{15.2cm}|>{\centering\arraybackslash}m{1.7cm}|}
\caption{Параметры ОНМНП(ОНМОП)\hfill\vspace{-0.5\baselineskip}}\label{tnznp:tbl1}\\ 
\hhline{|-|-|}
\rowcolor{gray!30}
Наименование параметра & Значение  \\ 
\hhline{|-|-|}
\endfirsthead

\caption*{\hspace{3pt}\emph{Продолжение таблицы \ref{tnznp:tbl1}\hfill\vspace{-0.5\baselineskip}}} \\
\hhline{|-|-|}
\rowcolor{gray!30}
Наименование параметра & Значение  \\ 
\hhline{|-|-|}
\endhead
\hhline{|-|-|}

\endfoot
\endlastfoot

\raggedright Время срабатывания ОНМ при утроенных по отношению к порогам срабатывания значениях тока и напряжения, мс, не более & \centering\arraybackslash 35  \\  \hline
\raggedright Время возврата ОНМ при одновременном сбросе входных величин тока и напряжения от номинальных значений до нуля, мс, не более  & \centering\arraybackslash 25  \\  \hline
\raggedright Средняя основная погрешность порогов срабатывания ОНМ по току и напряжению, \% уставки, не более   & \centering\arraybackslash $\pm$5  \\  \hline
\raggedright Средняя основная абсолютная погрешность ОНМ по углу максимальной чувствительности, эл. град., не более & \centering\arraybackslash $\pm$5 \\  \hline
\raggedright Дополнительная погрешность параметров срабатывания от изменения температуры окружающего воздуха в рабочем диапазоне, \% от среднего значения параметров, измеренных при температуре $(20\pm5)^\circ$C, не более & \centering\arraybackslash $\pm$5 \\  \hline
\raggedright Минимальная угловая ширина зоны срабатывания ОНМ, эл. град, не менее & \centering\arraybackslash 165 \\ \hline
\raggedright Коэффициент возврата ОНМ НП по напряжению и току, не менее  & \centering\arraybackslash 0,9\\  \hline

\end{longtable}
\setlength{\extrarowheight}{0.0cm} % возврат параметра высота ячеек в таблице
\normalsize

\item
На рисунке \ref{pict_tnznp:ONMNP} рисунке представлены характеристики срабатывания ОНМНП в 1СШ и ОНМНП в 2СШ. Аналогичные
характеристики для ОНМОП в 1СШ и ОНМОП в 2СШ.

\medskip

\begin{figure}[H]
\centering
\begin{subfigure}{0.5\textwidth}
    \centering
    \includegraphics[width=\textwidth, trim=0 0mm 125mm 0, clip]{\fbpath/661.0121 ТНЗНП/img/ONMNP.pdf}
	    \caption{a) ОНМНП в 1СШ}
    \label{pict_tnznp:ONMNP_a}
\end{subfigure}
%\hfill
\begin{subfigure}{0.49\textwidth}
    \centering
    \includegraphics[width=\textwidth, trim=125mm 0mm 0 0, clip]{\fbpath/661.0121 ТНЗНП/img/ONMNP.pdf}
    \caption{б) ОНМНП в 2СШ}
    \label{pict_tnznp:ONMNP_b}
\end{subfigure}

\caption{Характеристики срабатывания ОНМНП}
\label{pict_tnznp:ONMNP}
\end{figure}

\medskip

\item
Для выбора режима направленности по мощности нулевой последовательности для каждой ступени
ТНЗНП предусмотрен программный переключатель \textbf{«Направл. ТНЗНП-1(2…4)»}, имеющий три положения:
\begin{enumerate}[label=\arabic*)]
    \item \textit{«Ненапр»} -- ступень выполнена ненаправленной;
    \item \textit{«В 1СШ»} -- ступень выполнена прямонаправленной в сторону первой системы шин;
    \item \textit{«В 2СШ»} -- ступень выполнена обратнонаправленной в сторону второй системы шин.
\end{enumerate}

\item
Кроме того, можно установить различные алгоритмы контроля направленности с помощью программных
переключателей \textbf{«Контр. направл. ТНЗНП-1(2…4)»}, имеющих три положения:
\begin{enumerate}[label=\arabic*)]
    \item \textit{«Разр»} --  прямая направленность обеспечивается срабатыванием ОНМНП в 1СШ, обратная -- срабатыванием ОНМНП в 2СШ;
    \item \textit{«Блок»} -- прямая направленность обеспечивается отсутствием срабатывания ОНМНП в 2СШ, обратная -- отсутствием срабатывания ОНМНП в 1СШ;
    \item \textit{«Разр или блок»} -- прямая направленность обеспечивается срабатыванием ОНМНП в 1СШ или отсутствием срабатывания ОНМНП в 2СШ, обратная -- срабатыванием ОНМНП в 2СШ или отсутствием срабатывания ОНМНП в 1СШ.
\end{enumerate}

\item
Положение переключателя \textit{«Разр или блок»} рекомендуется выставлять для ступеней ТНЗНП,
выполняющих функцию дальнего резервирования, так как при удаленных КЗ напряжение $3U_0$ может оказаться
недостаточным для срабатывания ОНМ из-за его уменьшения по мере удаления от точки КЗ.

\item
Уставки для выбора режима направленности по мощности обратной последовательности аналогичные.

\item
Тип ОНМ, используемый для определения направленности ступени ТНЗНП, выбирается программным
переключателем \textit{«Тип ОНМ ТНЗНП-1(2…4)»}, имеющим четыре положения:
\begin{enumerate}[label=\arabic*)]
    \item \textit{«НП»} -- для определения направленности ступени учитывается только ОНМНП;
    \item \textit{«ОП»} -- для определения направленности ступени учитывается только ОНМОП;
    \item \textit{«НП или ОП»} -- для определения направленности ступени учитывается один из органов, ОНМНП или ОНМОП;
    \item \textit{«НП и ОП»} -- для определения направленности ступени одновременно учитываются оба органа, ОНМНП и ОНМОП.
\end{enumerate}

\item
Имеется возможность задать различные режимы работы ступеней ТНЗНП при возникновении
неисправности в цепях напряжения. Вариант действия логики ступени ТНЗНП при
срабатывании БНН задается программным переключателем \textbf{«БНН ТНЗНП-1(2…4)»}, имеющим три положения:
\begin{enumerate}[label=\arabic*)]
\item \textit{«Выведено»} -- режим работы ступени не зависит от БНН;
\item \textit{«ОНМ»} -- ступень переводится в ненаправленный режим работы;
\item \textit{«ТНЗНП»} --  при срабатывании БНН ступень выводится из работы Однако, если ступень задана ненаправленной, то она останется в работе независимо от срабатывания БНН. Ступень также не будет выводиться из работы, если направленность автоматически вывелась при срабатывании ТНЗНП.
\end{enumerate}

\item
Предусмотрен автоматический вывод направленности всех ступеней ТНЗНП при срабатывании ТНЗНП,
это обеспечивает устойчивое срабатывание ТНЗНП при неполнофазном отключении, что необходимо для
надежного действия УРОВ. Данный режим вводится с помощью программного переключателя \textbf{«Вывод
направл при ср ТНЗНП»}.

\item
Для обеспечения несрабатывания ТНЗНП от бросков намагничивающего тока предусмотрена
блокировка ступеней ТНЗНП по содержанию второй гармоники в токе нулевой последовательности.
Блокировка основана на контроле отношения действующего значения второй гармоники тока нулевой
последовательности к основной гармонике. Ступени ТНЗНП блокируется при превышении данного
соотношения пороговой величины, заданной уставкой \textbf{«3I02г/3I01г ТНЗНП»}. Наличие функции блокировки от
бросков намагничивающего тока задается для каждой ступени отдельно с помощью программных ключей
\textbf{«Режим БНТ ТНЗНП-1(2…4)»}. Для предотвращения излишней блокировки при больших токах КЗ,
сопровождающихся насыщением ТТ, блокировка выводится из действия при превышении тока 3I0 основной
гармоники выше уставки \textbf{«3I0блок ОВ БНТ ТНЗНП»}.

\item
Оперативный ввод/вывод чувствительных ступеней ТНЗНП производится с помощью команды \textit{«Вывод
ЧС ТНЗНП»}, а выбор выводимых ступеней ТНЗНП определяется с помощью программных ключей \textbf{«Режим
БЧС ТНЗНП-1(2…4)»}.

\item
Оперативный ввод/вывод ТНЗНП производится с помощью БУ (блоков управления функциями) \textit{«Вывод ТНЗНП»}. 

\item
В составе логики токовой направленной защиты нулевой последовательности предусмотрено два
режима оперативного ускорения ТНЗНП..

\vspace{2mm}
\begin{figure}[H]
\centering
\includegraphics[width=0.9\textwidth,height=0.9\textheight,keepaspectratio]{\fbpath/661.0121 ТНЗНП/img/OA_TNZNP_without_T.pdf}
\captionof{figure}{Функциональная схема ОУ ТНЗНП без выдержки времени}\label{pict_tnznp:OA_TNZNP_without_T}
\end{figure}

\item
Оперативное ускорение ТНЗНП без выдержки времени предназначено для оперативного ускорения
ступени ТНЗНП при выполнении переключений в сети. Ввод в работу ОУ-бв ТНЗНП и выбор ускоряемой
ступени задается программным переключателем \textbf{«ОУ-бв ТНЗНП»}. Предусмотрена возможность вывода
направленности ускоряемой ступени с помощью программного ключа \textbf{«Направл ТНЗНП при ОУ-бв»}. Действие
ускорения осуществляется с регулируемой независимой выдержкой времени \textbf{«Тср ОУ-бв ТНЗНП»}. При
выводе направленности ускоряемой ступени ее блокировка от БНТ осуществляется с помощью программного
ключа \textbf{«Режим БНТ ОУ-бв ТНЗНП ненапр»}.

\item
Оперативный ввод ОУ-бв ТНЗНП производится с помощью команды управления \textit{«Ввод ОУ-бв ТНЗНП»}.

\vspace{2mm}
\begin{figure}[H]
\centering
\includegraphics[width=0.9\textwidth,height=0.9\textheight,keepaspectratio]{\fbpath/661.0121 ТНЗНП/img/OA_TNZNP.pdf}
\captionof{figure}{Функциональная схема ОУ ТНЗНП}\label{desc_tnznp:OA_TNZNP}
\end{figure}

\item
Оперативное ускорение ТНЗНП предназначено для оперативного ускорения ступени ТНЗНП при выводе
ДЗШ. Ввод в работу ОУ ТНЗНП и выбор ускоряемой ступени задается программным переключателем \textbf{«ОУ
ТНЗНП»}. Предусмотрена возможность вывода направленности ускоряемой ступени с помощью программного
ключа \textbf{«Направл ТНЗНП при ОУ»}. Действие ускорения осуществляется с регулируемой независимой
выдержкой времени \textbf{«Тср ОУ ТНЗНП»}. При выводе направленности ускоряемой ступени ее блокировка от БНТ
осуществляется с помощью программного ключа \textbf{«Режим БНТ ОУ ТНЗНП ненапр»}.

\item
Оперативный ввод ОУ ТНЗНП производится с помощью команды управления \textit{«Ввод ОУ ТНЗНП»}.

\item
В составе логики токовой направленной защиты нулевой последовательности предусмотрен пуск
автоматического ускорения ступеней ТНЗНП.

\vspace{2mm}
\begin{figure}[H]
\centering
\includegraphics[width=0.9\textwidth,height=0.9\textheight,keepaspectratio]{\fbpath/661.0121 ТНЗНП/img/AA_TNZNP.pdf}
\captionof{figure}{Функциональная схема пуска АУ ТНЗНП}\label{pict_tnznp:AA_TNZNP}
\end{figure}

\item
Ввод в работу пуска АУ ТНЗНП и выбор ускоряемой ступени задается программным переключателем
\textbf{«АУ ТНЗНП»}. Предусмотрена возможность вывода направленности ускоряемой ступени с помощью
программного ключа \textbf{«Направл ТНЗНП при АУ»}. При выводе направленности ускоряемой ступени ее
блокировка от БНТ осуществляется с помощью программного ключа \textbf{«Режим БНТ АУ ТНЗНП ненапр»}.

\small
\begin{longtable}{|>{\centering\arraybackslash}m{5.3cm}|>{\centering\arraybackslash}m{3.3cm}|>{\centering\arraybackslash}m{4.2cm}|>{\centering\arraybackslash}m{1.8cm}|>{\centering\arraybackslash}m{1cm}|}
\caption{Перечень уставочных параметров ТНЗНП\hfill\vspace{-0.5\baselineskip}}\label{tnznp:tbl2}\\ 
\hline
\rowcolor{gray!30}
Параметр & Условное обозначение на схеме & Значение/ Диапазон & Единица измерения & Шаг \\ 
\hline
\endfirsthead
\caption*{\hspace{3pt}\emph{Продолжение таблицы \ref{tnznp:tbl2}\hfill\vspace{-0.5\baselineskip}}} \\
\hline
\rowcolor{gray!30}
Параметр & Условное обозначение на схеме & Значение/ Диапазон & Единица измерения & Шаг \\ 
%\hline
\endhead

%===m>ТНЗНП|-
%===m

\end{longtable}
\normalsize

\end{enumerate}